% Experimental setup (H/w and S/w requirements). Evaluation metrics (including
% equations and detailed metric descriptions). Experiments planned

\subsection{Experimental Setup} \label{subsection:experimental_setup}

The project is developed in Python and utilizes libraries mentioned in table
\ref{table:software_dependencies} and implemented on a machine with the
specifications shown in the table \ref{table:hardware_specs}.

\begin{table}[h!t]
  \caption{Software Dependencies}
  \begin{center}
    \begin{tabularx}{\linewidth}{rllY}
      \toprule
        & \textbf{Libraries} & \textbf{Versions} & \textbf{Description} \\
      \midrule
      1 & PyTorch            & 2.11.0+cu130      & Machine Learning Framework \\
      2 & Torchvision        & 0.26.0+cu130      & Computer Vision Library \\
      3 & OpenCV             & 4.13.0            & Computer Vision Library \\
      4 & Pillow             & 12.1.1            & Image Processing Library \\
      5 & Scikit-learn       & 1.8.0             & Machine Learning Library \\
      6 & NumPy              & 2.4.3             & Numerical Computing Library \\
      7 & SciPy              & 1.17.1            & Scientific Computing Library \\
      7 & Matplotlib         & 3.9.1             & Data Visualization Library \\
      7 & Seaborn            & 0.13.2            & Data Visualization Library \\
      \bottomrule
    \end{tabularx}
    \label{table:software_dependencies}
  \end{center}
\end{table}

\begin{table}[h!t]
  \caption{Hardware Specifications}
  \begin{center}
    \begin{tabularx}{\linewidth}{rlY}
      \toprule
        & \textbf{Hardware/Software} & \textbf{Specifications} \\
      \midrule
      1 & Operating System           & Windows 10/11 \\
      2 & Programming Language       & Python 3.11 \\
      3 & Memory                     & 32GB \\
      4 & CPU                        & AMD Ryzen 9 7940HX 2.4 GHz \\
      5 & GPU                        & NVIDIA GeForce RTX 4070 Laptop 8GB \\
      6 & Environment                & Miniconda \\
      \bottomrule
    \end{tabularx}
    \label{table:hardware_specs}
  \end{center}
\end{table}


\subsection{Evaluation Metrics} \label{subsection:evaluation_metrics}

The framework is evaluated on accuracy, area under the ROC curve,
precision, recall, and F1-score metrics. This concept has achieved an
accuracy of $92\%$ and an AUC of $0.98$ on the test set. The full detailed
results of the metrics will be discussed in section \ref{section:results}.

\subsubsection{Accuracy}
\begin{equation}
  \label{eqn:accuracy}
  \text{Accuracy} = \frac{T_P + T_N}{T_P + T_N + F_P + F_N}
\end{equation}

Where, $T_P$ stands for True Positives, $T_N$ for True Negatives, $F_P$ for
False Positives and $F_N$ for False Negatives.
Accuracy, as defined in Equation \ref{eqn:accuracy}, is the ratio of correctly
classified instances, i.e. the sum of true positives ($T_P$) and true negatives
($T_N$), to the total number of instances.
\\
\subsubsection{Area Under the ROC Curve}
The Receiver Operating Characteristic (ROC) Curve is a visualization of the
performance of a binary classifier across different thresholds.The Area Under
the ROC Curve (AUC) is a metric that summarizes the overall performance of the
classifier. An AUC of $1.0$ shows perfect classification, while an AUC of $0.5$
shows random guessing. The AUC is calculated by integrating the area under the
ROC curve, which plots the True Positive Rate ($T_PR$) against the False
Positive Rate ($F_PR$) at various threshold settings. A higher AUC value
indicates better model performance in distinguishing between the positive and
negative classes. \\
\subsubsection{Precision}
Precision is the ratio of true positives to the sum of true and false
positives, measuring the accuracy of the positive predictions made by
the model, as shown in the equation
\begin{equation}
  \label{eqn:precision}
  \text{Precision} = \frac{T_P}{T_P + F_P}
\end{equation}

Where, $T_P$ stands for True Positives and $F_P$ for False Positives.
A high precision is an indicator for a model that makes fewer false
positive errors, meaning that when it predicts a positive class, it
is likely to be correct.
\\
\subsubsection{Recall}
Recall is that machine learning metric, evaluating a model's ability
to find all relevant cases in a dataset, mathematically calculated as
the ratio of true positives to the total number of actual positives,
which is the sum of true positives and false negatives in the
predictions. Recall is calculated using equation \ref{eqn:recall}
\begin{equation}
  \label{eqn:recall}
  \text{Recall} = \frac{T_P}{T_P + F_N}
\end{equation}

Where, $T_P$ stands for True Positives and $F_N$ for False Negatives.
\\
\subsubsection{F1-Score}
That evaluation metric in machine learning defined as the harmonic
mean of precision and recall is the F1 Score, combining both into a
balanced score ranging from 0 to 1; 0 being the worst and 1 being the
best. A higher F1 Score indicates a better performance. The formula
for the F1 Score is mentioned in equation \ref{eqn:f1_score}
\begin{equation}
  \label{eqn:f1_score}
  \begin{split}
    \text{F1} &= \frac{2 \times (Precision \times Recall)}{Precision
    + Recall} \\
    &= \frac{2 \times T_P}{(2 \times T_P) + F_P + F_N}
  \end{split}
\end{equation}


\subsubsection{Precision-Recall Curve}
The Precision-Recall (PR) curve is a graphical evaluation metric
which plots precision against recall across variable classification
thresholds. PR curves are particularly used for imbalance datasets
thanks to their focus on positive-class prediction quality over ROC
curves. A better classification performance occurs when the PR has a
larger area underneath the curve.

\subsubsection{Macro Averaged Metrics}
Macro averaging computes the metric independently for each class and
then averages them equally, treating all classes equally regardless
of the dataset size. It is represented by equation \ref{eqn:macro_avg}
\begin{equation}
  \label{eqn:macro_avg}
  \text{Macro Average} = \frac{1}{N} \sum_{i=1}^{N} M_i
\end{equation}

where $M_i$ is the metric value for each class.

\subsubsection{Weighted Averaged Metrics}
Weighted average deals with the average metric while taking into
account the number of samples in each case and assigning varying
levels of importance to data points or models, denoted by equation
\ref{eqn:weighted_avg}
\begin{equation}
  \label{eqn:weighted_avg}
  \text{Weighted Average} = \frac{\sum_{i=1}^{N} w_i M_i}{\sum_{i=1}^{N} w_i}
\end{equation}

where $w_i$ is the class support and $M_i$ represents the metric value for
each class.
